\documentclass[a4paper]{article}

\usepackage{import}
\usepackage[utf8]{inputenc}
\usepackage[T1]{fontenc}
\usepackage{textcomp}
\usepackage[italian]{babel}
\usepackage{amsmath, amssymb}
\usepackage{booktabs,xltabular}
\usepackage{amsfonts}
\usepackage{subcaption}
\usepackage{amsthm}
\usepackage{cancel}
\usepackage{mdframed}
\usepackage{makecell}
\usepackage{float}
\usepackage[dvipsnames]{xcolor}
\usepackage{listings}
\usepackage{gensymb}
\usepackage{graphicx}
\usepackage[table]{xcolor}
\usepackage{booktabs}
\usepackage{tabularx}
% \usepackage{bodeplot}
\usepackage{tikz}
\usetikzlibrary{shapes, arrows, automata, petri, decorations.markings, decorations.pathreplacing, positioning, calc, quotes}
\usepackage{circuitikz}
% \usepackage[label=corner]{karnaugh-map}
\graphicspath{{./figures/}}

% Prevent overfull hbox warnings
\emergencystretch=1.5em

% Set default font to sans-serif
\renewcommand{\familydefault}{\sfdefault} 
% \usepackage{eulervm}

\usepackage{forest}

\usepackage{mathtools}
\DeclarePairedDelimiter\ceil{\lceil}{\rceil}
\DeclarePairedDelimiter\floor{\lfloor}{\rfloor}

% \usepackage{ntheorem}

\usepackage{import}
\usepackage{pdfpages}
% \usepackage{transparent}
\usepackage{xcolor}

\usepackage{hyperref}
\hypersetup{
    colorlinks=false,
}

% Code blocks
\definecolor{codegreen}{rgb}{0,0.6,0}
\definecolor{codegray}{rgb}{0.5,0.5,0.5}
\definecolor{codepurple}{rgb}{0.58,0,0.82}
\definecolor{backcolour}{rgb}{0.95,0.95,0.95}

\lstdefinestyle{mystyle}{
	backgroundcolor=\color{backcolour},
	commentstyle=\color{codegreen},
	keywordstyle=\color{magenta},
	numberstyle=\tiny\color{codegray},
	stringstyle=\color{codepurple},
	basicstyle=\ttfamily\footnotesize,
	breakatwhitespace=false,
	breaklines=true,
	captionpos=b,
	keepspaces=true,
	numbers=left,
	numbersep=5pt,
	showspaces=false,
	showstringspaces=false,
	showtabs=false,
	tabsize=2
}

\definecolor{head}{HTML}{1F3A5F}
\definecolor{rowgray}{HTML}{EEF2F7}

\renewcommand{\arraystretch}{1.5}
\newcolumntype{Y}{>{\raggedright\arraybackslash}X}

\lstdefinestyle{terminal}{
    backgroundcolor=\color{black},
    basicstyle=\ttfamily\color{green},
    frame=single,
    breaklines=true,
    showstringspaces=false
}

\lstset{style=mystyle}

\usepackage{color}
\usepackage{import}
\usepackage{pdfpages}
% \usepackage{transparent}
\usepackage{xcolor}

% Example frame
\theoremstyle{definition}
\newmdtheoremenv[%
	linecolor=gray,leftmargin=0,%
	rightmargin=0,
	innertopmargin=8pt,%
	innerbottommargin=8pt,
	ntheorem]{example}{Esempio}[section]

% Important definition frame
\theoremstyle{definition}
\newmdtheoremenv[%
	linecolor=gray,leftmargin=0,%
	rightmargin=0,
	backgroundcolor=gray!40,%
	innertopmargin=8pt,%
	innerbottommargin=8pt,
	ntheorem]{definition}{Definizione}[section]

% Exercise frame
\theoremstyle{definition}
\newmdtheoremenv[%
	linecolor=gray,leftmargin=0,%
	rightmargin=0,
	innertopmargin=8pt,%
	innerbottommargin=8pt,
	ntheorem]{exercise}{Esercizio}[section]

% Theorem frame
\theoremstyle{definition}
\newmdtheoremenv[%
  linecolor=gray,leftmargin=0,%
  rightmargin=0,
  innertopmargin=8pt,%
  innerbottommargin=8pt,
  ntheorem]{theorem}{Teorema}[section]

\theoremstyle{definition}
\newmdtheoremenv[%
  linecolor=white,leftmargin=0,%
  rightmargin=0,
  innertopmargin=8pt,%
  innerbottommargin=8pt,
  ntheorem]{define}{Definizione utile}[section]

% figure support
\usepackage{import}
\usepackage{xifthen}
\pdfminorversion=7
\usepackage{pdfpages}
% \usepackage{transparent}
\newcommand{\incfig}[1]{%
	\def\svgwidth{\columnwidth}
	\import{./figures/}{#1.pdf_tex}
}

% FSM tikz
\tikzset{
    place/.style={
        circle,
        thick,
        draw=black,
        minimum size=6mm,
    },
        state/.style={
        circle,
        thick,
        draw=black,
        fill=white,
        minimum size=6mm,
    },
}

\pdfsuppresswarningpagegroup=1

\usepackage{pgfplots}
\pgfplotsset{compat=1.18,width=10cm}

% Save plots as pdf and reuse them without compiling every time
\usetikzlibrary{external}
\tikzexternalize[prefix=figures/tikz/, optimize=false]


\begin{document}
\begin{titlepage}
	\begin{center}
		\vspace*{1cm}

		\Huge
		\textbf{Natural Language Processing}

		\vspace{0.5cm}
		\LARGE
		UniVR - Dipartimento di Informatica

		\vspace{1.5cm}

		\textbf{Tommaso Toffali}

		\vfill


		\vspace{0.8cm}


		1° Semestre 2026/2027
    \end{center}
\end{titlepage}

\tableofcontents
\pagebreak

\section{Introduction}
Language is a structured system of conventional signs used by agents to represent,
express and communicate meanings.

\subsection{How did language emerge?}
Two question must be distinguished:
\begin{itemize}
    \item \textbf{fThe evolutionary origin of the human capacity for language:} how biological, cognitive and 
        social evolution made language possibile.
    \item \textbf{The historical development of particular language:} how language changes, differentiate, interact
        and give rise to language families
\end{itemize}
The origin of language itself is not directly documented. Spoken and signed languages existed long before
writing, so their early history must be reconstructured indirectly using evidence from:
\begin{itemize}
    \item human biology and cognition
    \item archaeology
    \item comparative linguistic
    \item animal comunication
\end{itemize}

\subsection{What is language?}
Language is a structured system of conventional signs used by agents to represent, express and communicate meanings.
\begin{itemize}
    \item Spoken sound
    \item Signed gesture
    \item Written symbols
    \item Tactile signs, such as Braille
    \item Haptic signs, conveyed through touch or movement
\end{itemize}

\subsection{Who does use language?}
\begin{itemize}
    \item \textbf{Human agents}
    \item \textbf{Non-human animal agents}
        \begin{itemize}
            \item some non-human animals use structured communication system, although their classification
                as languages is disputed
        \end{itemize}
    \item \textbf{Artificial agents}
        \begin{itemize}
            \item artificial agents process and generate linguistic expressions, while the nature of their understanding remains an open theoretical question.
        \end{itemize}
\end{itemize}
\subsection{Level of linguistic representation}

\rowcolors{2}{rowgray}{white}
\begin{tabularx}{\textwidth}{>{\bfseries}p{3.6cm} Y Y}
    \rowcolor{head}
    \color{white}\bfseries Level &
    \color{white}\bfseries Principal units &
    \color{white}\bfseries Central question \\
    \toprule
    Signal & Sounds, gestures, written or tactile signs & How is language physically conveyed? \\
    Phonological or graphemic & Phonemes, syllables, graphemes & How are elementary signs organized and distinguished? \\
    Morphological & Morphemes and word forms & How are meaningful components combined into words? \\
    Lexical & Words, lemmas and senses & What words exist, and what can they mean? \\
    Syntactic & Phrases, clauses and sentences & How are words combined into structured expressions? \\
    Semantic & Concepts, entities, predicates, arguments and propositions & What does the expression mean? \\
    Discourse & Sentences, utterances and discourse relations & How does an expression connect with the surrounding text or conversation? \\
    Pragmatic & Speakers, hearers, intentions and context & What does an agent want to obtain by using the expression in this situation? \\
    \bottomrule
\end{tabularx}

\subsection{Typology based on writing system}
Languages should be distinguished based on their features. We may start by their writing system to step 
onto four further basic features:
\begin{itemize}
    \item Morphology
    \item Syntax
    \item Semantics
    \item Pragmatics
\end{itemize}

\subsection{IPA}
The international phonetic alphabet (IPA) is a standardized system for representing the sounds of spoken language.
Ideally, each symbol identifies a particular speech sound independently of the ordinary spelling conventions of a laguage.
\vspace{1em}
\noindent
\newline
We can present three languages as a progression in the relationship between \textbf{written form and pronounciation}.
This is not a progression from a \textit{simple language} to a \textit{complex language}, it concerns the complexity
of their orthographic systems and form-sound correspondences. 

\subsection{Ambiguity}
Ambiguity occurs when the same linguistic expression permits two or more distinct analyses or interpretations.
\end{document}